\subsection{Relaciones entre imágenes directas e imágenes inversas}
\label{subsection:0.4.4-fr}

\begin{env}[4.4.1]
\label{0.4.4.1-fr}
Con las hipótesis y la notación de (4.2.1), sea $\mathcal{G}$ un
$\mathcal{B}$-Módulo. Por definición, un homomorfismo
$u:\mathcal{G}\to\Psi_*(\mathcal{F})$ de $\mathcal{B}$-Módulos se llama
también \emph{$\Psi$-morfismo de $\mathcal{G}$ en $\mathcal{F}$}, o
simplemente \emph{homomorfismo de $\mathcal{G}$ en $\mathcal{F}$}, y se
escribe $u:\mathcal{G}\to\mathcal{F}$ cuando no hay lugar a confusión.
Dar tal homomorfismo equivale a dar, para todo par $(U,V)$ en el que
$U$ es un abierto de $X$ y $V$ un abierto de $Y$ con
$\psi(U)\subset V$, un \emph{homomorfismo}
\[
  u_{U,V}:\Gamma(V,\mathcal{G})\longrightarrow\Gamma(U,\mathcal{F})
\]
\emph{de $\Gamma(V,\mathcal{B})$-módulos}, donde $\Gamma(U,\mathcal{F})$
se considera como $\Gamma(V,\mathcal{B})$-módulo mediante el homomorfismo
de anillos
$\theta_{U,V}:\Gamma(V,\mathcal{B})\to\Gamma(U,\mathcal{A})$. Además,
los $u_{U,V}$ deben hacer conmutativos los diagramas (3.5.1.1). Por lo demás,
para definir $u$ basta dar los $u_{U,V}$ cuando $U$ (resp. $V$) recorre
una base $\mathfrak{B}$ (resp. $\mathfrak{B}'$) de la topología de $X$
(resp. $Y$) y verificar la conmutatividad de (3.5.1.1) para estas
restricciones.
\end{env}
